\documentclass[10pt, letterpaper]{article}

\usepackage[
    ignoreheadfoot,
    top=2cm,
    bottom=2cm,
    left=2cm,
    right=2cm,
    footskip=0.8cm,
]{geometry}
\usepackage{fontawesome5}
\usepackage{hyperref}
\hypersetup{colorlinks=true, urlcolor=black}
\usepackage[french]{babel}
\usepackage[T1]{fontenc}
\usepackage{lmodern}

\pagestyle{empty}
\setlength{\parindent}{0pt}
\setlength{\parskip}{0.5em}

\begin{document}

\begin{center}
    \textbf{\LARGE Candidature -- Graduate Quantitative Analyst (Risk \& Model Validation)} \\[0.2cm]
    \textbf{\large Charly-Romy TANGA} \\[0.2cm]
    \href{mailto:charlyromytanga.cytech@gmail.com}{\faEnvelope\ charlyromytanga.cytech@gmail.com}
    \quad$\cdot$\quad
    \faPhone\ +33 6 15 42 25 74
    \quad$\cdot$\quad
    \href{https://www.linkedin.com/in/charly-romy-tanga}{\faLinkedin\ charly-romy.tanga}
    \quad$\cdot$\quad
    \href{https://github.com/charlyromytanga}{\faGithub\ charlyromytanga}
\end{center}

Madame, Monsieur,

Actuellement en dernière année d'école d'ingénieur en mathématiques appliquées et informatique, je souhaite débuter ma carrière en conseil quantitatif, gestion des risques ou validation de modèles.

Mon alternance m'a permis d'évoluer dans un environnement exigeant où la qualité des données, la rigueur méthodologique et la collaboration avec les équipes métiers sont essentielles. Cette expérience m'a donné le goût des environnements structurés, de la résolution de problématiques complexes et du travail en équipe.

En parallèle, je développe des projets en finance quantitative autour de la simulation Monte-Carlo, du calcul de Value-at-Risk et du pricing d'options, afin de renforcer mes compétences en modélisation et méthodes numériques.

Motivé, rigoureux et curieux, je souhaite rejoindre un cabinet où je pourrai contribuer à des missions de validation de modèles, de gestion des risques ou de transformation data des institutions financières.

Je serais ravi d'échanger avec vous lors d'un entretien.

Cordialement,\\
Charly-Romy TANGA

\end{document}
